% ============================================================
\section{Métricas e métodos de análise}\label{sec:metodos}
% ============================================================

% ------------------------------------------------------------
\subsection{Métricas estruturais}\label{sec:met-estruturais}

Em todas as redes da configuração principal calculamos as métricas exigidas pela disciplina,
mais reciprocidade e \emph{hubs}. Seja $G=(V,E)$ a rede não direcionada por união, $n=|V|$,
$m=|E|$ e $k_i$ o grau do vértice $i$.
\begin{description}
    \item[Tamanho e grau.] $n$, $m$, grau médio $\langle k\rangle=2m/n$ e densidade
    $\rho=2m/[n(n-1)]$. Nas redes k-NN, $m$ fica entre $nk/2$ (todas as escolhas recíprocas) e
    $nk$ (nenhuma recíproca), então $\langle k\rangle\in[k,2k]$ e a densidade é muito baixa
    por construção; o interesse está em como $m$ varia entre camadas, que reflete a
    reciprocidade.
    \item[Distribuição de graus.] Na rede não direcionada, e de entrada e de saída na
    direcionada. Na variante (a), o grau de saída é $k$ para todos, e toda a informação está no
    grau de entrada $k^{\mathrm{in}}_i=|\{j: i\in\N_j\}|$, cuja média também é $k$. Reportamos a
    CCDF $P(K^{\mathrm{in}}\ge x)$ e a assimetria (\emph{skewness}) de $k^{\mathrm{in}}$, a medida
    usual de \emph{hubness}~\cite{radovanovic2010}: assimetria alta significa poucos vértices
    escolhidos por muitos.
    \item[Reciprocidade.] Fração das arestas direcionadas $i\to j$ para as quais $j\to i$
    também existe.
    \item[Clusterização.] Global (transitividade),
    $C=3\times(\text{triângulos})/(\text{trios conexos})$, e local média,
    \begin{equation}
        \bar C=\frac{1}{|V_{\ge2}|}\sum_{i\in V_{\ge2}} C_i,\qquad C_i=\frac{2\,e_i}{k_i(k_i-1)},
    \end{equation}
    onde $e_i$ é o número de arestas entre os vizinhos de $i$. $C_i$ não é definido quando
    $k_i<2$, e esses vértices ficam fora da média: $V_{\ge2}$ é o conjunto dos vértices de
    grau $\ge 2$ (na rede por união $k_i\ge k$, então $V_{\ge2}=V$).
    \item[Distâncias.] Distância média $\langle\ell\rangle$ entre pares do maior componente
    conexo, diâmetro e histograma das distâncias, por busca em largura exata a partir de todos
    os vértices.
    \item[Componentes.] Número de componentes conexos, fração de vértices no maior componente
    e distribuição dos tamanhos, para toda representação e nas versões por união e mútua. A
    versão por união \emph{não} é necessariamente conexa: um tipo com $f_t>k$ manda todas as
    $k$ arestas de saída de cada ocorrência para outras ocorrências do próprio tipo ($s=1$) e
    vira um componente fechado se nenhum vértice de outro tipo escolher alguma delas; o mesmo
    vale, nas camadas contextuais, para um grupo de prefixo com mais de $k$ ocorrências
    (vetores idênticos) e, na rede do vocabulário, para um grupo de mais de $k$ tipos com
    \emph{embeddings} idênticos (há 58 grupos desses, o maior com 524 tipos;
    Seção~\ref{sec:red-vocab}). Espera-se, portanto, que a rede lexical por união tenha vários
    componentes pequenos formados justamente pelos tipos frequentes do desenho (alvos,
    controles e multitemáticos); numa simulação com o perfil de frequências da amostra híbrida,
    $k=10$ e vetores aleatórios, foram 23 a 24 componentes, com o maior reunindo cerca de 96\%
    dos vértices. Como distância média e diâmetro são calculados no maior componente, na rede
    lexical eles excluem esses tipos, e a comparação de $\langle\ell\rangle$ entre
    representações deve ser lida junto com a fração do maior componente, que a tabela de
    métricas da rede por união informa.
\end{description}

As métricas são calculadas com o igraph~\cite{csardi2006}, cuja busca em largura roda em C:
as distâncias exatas de uma rede de 15 mil vértices levam segundos, contra 20 a 25 minutos por
rede no NetworkX do piloto. O NetworkX fica como referência nos testes (as duas bibliotecas
devem dar os mesmos valores em grafos pequenos). Na rede do vocabulário (151,6 mil vértices),
as distâncias exatas são calculadas só na configuração principal (estimativa: de dezenas de
minutos a cerca de uma hora); nas variantes de robustez, a distância média é estimada por busca
em largura a partir de 500 fontes sorteadas, e o diâmetro, por um limite inferior de varredura
dupla (a partir de um vértice, ir ao mais distante e repetir).

% ------------------------------------------------------------
\subsection{Persistência da vizinhança: Jaccard}\label{sec:met-jaccard}

A reorganização da vizinhança de uma ocorrência entre as redes $a$ e $b$ (por exemplo, lexical
e bloco 1) é medida pela similaridade de Jaccard
\begin{equation}
    J_i(a,b)=\frac{|\N_i^{(a)}\cap\N_i^{(b)}|}{|\N_i^{(a)}\cup\N_i^{(b)}|}.
\end{equation}
Com vizinhanças de tamanho fixo $k$ (variante (a)) e $c=|\N_i^{(a)}\cap\N_i^{(b)}|$ vizinhos em
comum, $J_i=c/(2k-c)$: $J=1$ quando nada muda e $J=0$ quando nenhum vizinho se mantém. Na
variante (b), os conjuntos têm tamanhos diferentes e a fórmula geral vale; na variante (c), a
mesma fórmula é aplicada aos conjuntos de tipos $T_i$. Chamamos de \emph{mudança} o complemento
$M_i=1-J_i$. Com 20 sementes na rede lexical, $J_i(\mathrm{lex},\ell)$ é calculado para cada
semente e resumido por média e desvio.

% ------------------------------------------------------------
\subsection{Dominância lexical e seu teto}\label{sec:met-dominancia}

A dominância lexical é a fração de vizinhos do mesmo tipo,
\begin{equation}
    D_i=\frac{|\{\,j\in\N_i : t_j=t_i\,\}|}{|\N_i|},
\end{equation}
com $|\N_i|=k$ na variante (a) e o tamanho variável na variante (b). Na rede lexical com a
variante (a), as outras ocorrências do tipo são as mais similares ($s=1$), então
\begin{equation}
    D_i^{\mathrm{lex}}=D^{\max}_i=\frac{\min(f_{t_i}-1,\,k)}{k}
\end{equation}
exatamente (um teste verifica isso). Esse teto depende de $f_t$: uma ocorrência de um tipo com
$f_t=4$ nunca passa de $D=0{,}3$ com $k=10$. Para comparar faixas de frequência, reportamos
também a dominância relativa ao teto, $\tilde D_i=D_i/D^{\max}_i$ (definida para $f_t\ge2$ e
analisada para $f_t\ge3$). A queda $\Delta D$ entre duas camadas mede quanto a identidade do
token perde espaço na vizinhança.

% ------------------------------------------------------------
\subsection{Piso de ruído}\label{sec:met-ruido}

Na rede lexical, parte dos vizinhos é sorteada entre candidatos empatados. Duas construções com
sementes diferentes, $r\neq r'$, já dão $J_i(\mathrm{lex}_r,\mathrm{lex}_{r'})<1$ sem nenhuma
mudança de representação. O \textbf{piso de ruído} é essa mudança produzida só pelo desempate,
\begin{equation}
    M^{\text{ruído}}_i = 1-\frac{1}{m}\sum_{q=0}^{m-1}
        J_i\bigl(\mathrm{lex}_{2q},\mathrm{lex}_{2q+1}\bigr),\qquad m=\lfloor R/2\rfloor,
\end{equation}
isto é, a média sobre os pares \emph{disjuntos} de sementes $(0,1),(2,3),\dots$ (com $R=20$
sementes, 10 pares; com $R$ ímpar a última fica de fora). Usar pares disjuntos, em vez de
todos os $\binom{R}{2}$ pares, mantém as estimativas de cada par independentes entre si; o
valor esperado é o mesmo. O piso é
calculado por ocorrência e resumido por faixa de frequência (ele é maior para tipos frequentes,
em que o bloco empatado é grande). Uma mudança lexical $\to$ contextual só é interpretada como
reorganização se $M_i(\mathrm{lex},\ell)$ ficar acima do piso da faixa correspondente.

% ------------------------------------------------------------
\subsection{Jaccard por composição de tipos}\label{sec:met-jw}

Como complemento que não depende da semente, comparamos a \emph{composição por tipos} das
vizinhanças. Seja $n^{(a)}_{iu}$ o número de vizinhos de $i$ do tipo $u$ na rede $a$. O Jaccard
ponderado é
\begin{equation}
    J^{w}_i(a,b)=\frac{\sum_u \min\bigl(n^{(a)}_{iu},n^{(b)}_{iu}\bigr)}
                      {\sum_u \max\bigl(n^{(a)}_{iu},n^{(b)}_{iu}\bigr)}.
\end{equation}
Na rede lexical, o bloco empatado da fronteira costuma ser formado por ocorrências de um único
tipo (as do tipo seguinte no ranking), e então qualquer sorteio dá a mesma composição: $J^w$ é
invariante à semente sempre que o bloco da fronteira tem um só tipo (um teste verifica isso em
dados sintéticos). Isso vale para quase todas as linhas, mas não para todas: dois tipos com
\emph{embeddings} idênticos, ou com cossenos a menos de $\eps$ um do outro, dividem o mesmo
bloco, e nessas linhas a composição varia com a semente. A fração dessas exceções é
reportada junto com $J^w$. Ele responde a uma pergunta um pouco
diferente de $J$: a ocorrência continua perto dos mesmos \emph{tipos}, ainda que de outras
ocorrências deles?

% ------------------------------------------------------------
\subsection{Estratificação}\label{sec:met-estratos}

As distribuições de $J$, $M$, $D$ e $\tilde D$ são reportadas por faixa de frequência (na
amostra e no corpus; Seção~\ref{sec:faixas}), por categoria de token, por estrato da amostra
(núcleo, alvos, controles, multitema) e com controle pela faixa de posição. Na P2, além da
mudança média por transição, contamos a fração de vértices cuja maior mudança acontece em cada
transição (lexical $\to$ 1, $1\to18$, $18\to36$) e traçamos as curvas $J(\ell,\ell+1)$ e
$D(\ell)$ nos 36 blocos, a partir do \emph{memmap} de todas as camadas.

% ------------------------------------------------------------
\subsection{Comunidades}\label{sec:met-comunidades}

As comunidades são detectadas na rede por união com o algoritmo de
Leiden~\cite{traag2019}, que melhora o Louvain~\cite{blondel2008} garantindo comunidades
conexas, maximizando a modularidade~\cite{newman2006}
\begin{equation}
    Q=\frac1{2m}\sum_{ij}\Bigl(A_{ij}-\gamma\frac{k_ik_j}{2m}\Bigr)\,\delta(c_i,c_j),
\end{equation}
com resolução $\gamma=1$ na configuração principal e $\gamma\in\{0{,}5;\,2\}$ na robustez. No
igraph, o objetivo padrão do \cod{community\_leiden} é o CPM (\emph{constant Potts model}),
cuja resolução tem outro significado (com resolução 1, o CPM tende a produzir comunidades
unitárias); por isso o código sempre passa \cod{objective\_function="modularity"}. O Leiden roda
com 10 sementes (o igraph usa o gerador \cod{random} do Python, que é semeado a cada execução);
fica a partição de maior modularidade, e a \textbf{estabilidade} é o NMI médio entre as
partições das 10 execuções. O Louvain serve de checagem.

% ------------------------------------------------------------
\subsection{Concordância entre partições (P3)}\label{sec:met-concordancia}

As comunidades de cada camada são comparadas com partições conhecidas dos vértices: tipo de
token, tema, artigo, parágrafo, sentença, próximo token previsto, próximo token real e faixa de
posição (esta como controle de artefato). Três medidas são usadas:
\begin{description}
    \item[NMI.] Informação mútua normalizada, $\mathrm{NMI}(U,V)=2\,I(U;V)/[H(U)+H(V)]$, via
    \cod{igraph.\allowbreak compare\_communities}. O NMI cresce com o número de classes mesmo entre
    partições sem relação nenhuma~\cite{vinh2010}, o que é grave aqui: ``parágrafo'' tem cerca
    de 80 classes no núcleo e ``tipo'', milhares. Por isso cada NMI é acompanhado de uma
    \textbf{linha de base de permutação}: o NMI médio entre as comunidades e 20 permutações
    aleatórias do rótulo, que preservam o número e o tamanho das classes. Reportamos o NMI
    bruto e $\mathrm{NMI}-\mathrm{NMI}_{\mathrm{perm}}$.
    \item[ARI.] Índice de Rand ajustado~\cite{hubert1985}, já corrigido pelo acaso (zero em
    média para partições independentes).
    \item[Pureza.] $P=\frac1n\sum_c\max_\lambda |c\cap\lambda|$: a fração de vértices que
    pertencem ao rótulo majoritário da sua comunidade. É fácil de interpretar, mas favorece
    comunidades pequenas, então é lida junto com as outras duas.
\end{description}
A P3 é analisada \textbf{só no núcleo}, onde os parágrafos estão inteiros e artigo e parágrafo
são rótulos distintos (dois parágrafos por artigo); nos estratos esparsos, cada parágrafo
contribui com uma ou poucas ocorrências e os rótulos de contexto quase coincidem com o tipo. O
tema também é analisado na rede inteira.

% ------------------------------------------------------------
\subsection{Separação intra-tipo (P4)}\label{sec:met-p4}

Para cada tipo com $f_t\ge10$ (alvos, controles, multitema e palavras funcionais separados):
\begin{itemize}
    \item \textbf{número efetivo de comunidades} ocupadas pelas suas ocorrências,
    $e^{H_t}$ com $H_t=-\sum_c p_{tc}\ln p_{tc}$ e $p_{tc}$ a fração das ocorrências de $t$
    na comunidade $c$ (vale 1 quando todas estão juntas; na rede lexical tende a 1, porque
    ocorrências do mesmo tipo são idênticas e ligadas entre si);
    \item \textbf{NMI entre comunidade e tema} dentro do tipo: as ocorrências se separam
    \emph{segundo o tema}, ou só se espalham?
    \item \textbf{autossimilaridade}: cosseno médio entre as próprias ocorrências do tipo
    (1 na camada lexical)~\cite{ethayarajh2019};
    \item \textbf{diferença entre o cosseno médio dentro do mesmo tema e entre temas};
    \item \textbf{separação pelo sentido anotado} (cerca de 20 ocorrências por alvo, em
    \cod{senses.csv}), que controla o fato de o tema ser só um indicador aproximado do sentido;
    \item \textbf{redes ego} de 2 ou 3 alvos, com layout fixo entre camadas, para inspeção
    qualitativa.
\end{itemize}
A hipótese H4 é testada comparando a distribuição dessas medidas entre alvos e controles,
que têm a mesma alocação de temas e cotas (Seção~\ref{sec:amostra}).

% ------------------------------------------------------------
\subsection{Rede do vocabulário}\label{sec:met-vocab}

Na rede lexical de tipos, além das métricas estruturais: distribuição de graus e \emph{hubs}
(quais tokens são escolhidos por muitos: tokens pouco treinados? pedaços de bytes?),
comunidades comparadas com a classe de escrita (NMI e ARI), a fração de tokens que aparecem no
corpus em cada comunidade e a vizinhança das palavras-alvo no vocabulário inteiro.

% ------------------------------------------------------------
\subsection{Extensão opcional: vizinhos no vocabulário}\label{sec:met-lente}

Para cada ocorrência e representação (lexical, blocos 1, 18 e 36, e bloco 36 normalizado),
calculamos o ranking das linhas da matriz de \emph{embeddings} inteira pelo cosseno (em
\emph{float32} na GPU; aqui não há tratamento de empates, porque só a posição no ranking é
usada) e medimos, camada a camada, a posição do \emph{próprio} token e a do \emph{próximo}
token nesse ranking, e o Jaccard entre os $k=20$ vizinhos no vocabulário de camadas
consecutivas. A pergunta: em que camada a ocorrência deixa de parecer com o próprio token e
passa a parecer com o próximo? Nas camadas intermediárias, a comparação mistura espaços que não
foram treinados para se alinhar; só o bloco 36 normalizado é exatamente o que a matriz de saída
lê. Duas checagens validam a extensão: na camada lexical, o próprio token é o primeiro do
ranking; no bloco 36 normalizado, a linha de \emph{maior produto interno} (os \emph{logits}
recalculados a partir de $E$) coincide com o argmax dos \emph{logits} do próprio modelo em
praticamente todas as ocorrências (\concordanciaPredicao{} na última execução). Pelo cosseno, a
primeira linha pode diferir, porque as linhas de $E$ não têm todas a mesma norma.
